\BeleskeID{08-s055}
% element: 08-s055:e03
% element: 08-s055:notes
Izlaz je napon otpornika, pa se početna vrednost određuje razlikom ulaznog napona i napona kondenzatora. Ulaz neposredno posle skoka iznosi $U$. Napon kondenzatora ne može trenutno da se promeni, pa se njegova vrednost pre skoka prenosi kao početno stanje posle skoka.

Pretpostavlja se da su pre trenutka $t_0$ završeni svi prethodni prelazni procesi. Zato su struje kroz kondenzator i otpornik jednake nuli, a pad napona na otporniku iščezava. Napon kondenzatora tada je jednak ranijem ulaznom naponu $U_0$. Razlika $U-U_0$ zato se u prvom trenutku pojavljuje na otporniku i predstavlja početni izlazni napon.
